\documentclass[10pt,letterpaper]{article}
\usepackage[utf8]{inputenc}
\usepackage[spanish,es-tabla]{babel}
\usepackage{geometry}
\geometry{margin=0.68in, top=0.85in, bottom=0.78in, headheight=36pt, headsep=14pt}
\usepackage{amsmath,amssymb}
\usepackage{xcolor}
\usepackage{graphicx}
\usepackage{tabularx}
\usepackage{booktabs}
\usepackage{colortbl}
\usepackage{array}
\usepackage{enumitem}
\usepackage{fancyhdr}
\usepackage{titlesec}
\usepackage{lastpage}
\usepackage{setspace}
\usepackage{listings}

% Tipografía institucional Icesi (Sans-Serif Helvetica)
\usepackage{helvet}
\renewcommand{\familydefault}{\sfdefault}

% Interlineado aumentado un 30% (setstretch 1.38)
\setstretch{1.38}

% Altura de celdas en tablas
\renewcommand{\arraystretch}{1.25}

% Paleta de Colores Oficial de los Programas Icesi
\definecolor{icesiBlue}{HTML}{5454E9}       % Azul vibrante institucional
\definecolor{icesiOrange}{HTML}{F0681A}     % Naranja acento
\definecolor{icesiGreen}{HTML}{4CB978}      % Verde acento
\definecolor{icesiPurple}{HTML}{865CF0}     % Púrpura acento
\definecolor{icesiDark}{HTML}{1A1A1A}       % Texto principal oscuro
\definecolor{icesiMuted}{HTML}{6B7280}      % Gris metadatos
\definecolor{icesiBorder}{HTML}{D1D5DB}     % Líneas y bordes
\definecolor{icesiRowLine}{HTML}{E5E7EB}    % Líneas de tabla
\definecolor{icesiLight}{HTML}{F9FAFB}      % Fondo suave
\definecolor{codeBackground}{HTML}{F3F4F6}  % Fondo código

% Configuración de Listings
\lstdefinestyle{bashstyle}{
    backgroundcolor=\color{codeBackground},
    commentstyle=\color{icesiMuted}\itshape,
    keywordstyle=\color{icesiBlue}\bfseries,
    basicstyle=\ttfamily\scriptsize,
    breaklines=true,
    frame=single,
    rulecolor=\color{icesiBorder},
    language=bash,
    aboveskip=2pt,
    belowskip=2pt
}

\lstdefinestyle{codeStyle}{
    backgroundcolor=\color{codeBackground},
    commentstyle=\color{icesiMuted}\itshape,
    keywordstyle=\color{icesiBlue}\bfseries,
    stringstyle=\color{icesiOrange},
    basicstyle=\ttfamily\scriptsize,
    breaklines=true,
    frame=single,
    rulecolor=\color{icesiBorder},
    aboveskip=2pt,
    belowskip=2pt
}

% =========================================================================
% ENCABEZADO Y PIE DE PÁGINA ESTILO PROGRAMA DE CURSO ICESI
% =========================================================================
\pagestyle{fancy}
\fancyhf{}
\fancyhead[L]{%
  \begin{minipage}[c]{2.8cm}
    \includegraphics[height=0.58cm]{logo-icesi.png}
  \end{minipage}%
  \hspace{0.15cm}%
  {\color{icesiBorder}\vrule width 0.8pt height 0.60cm}%
  \hspace{0.2cm}%
  \begin{minipage}[c]{10.5cm}
    \fontsize{7.5}{9.5}\selectfont\textbf{\color{icesiDark}UNIVERSIDAD ICESI} \\
    \fontsize{7}{9}\selectfont\textbf{\color{icesiDark}PROGRAMA DE CURSO / GUÍA PRÁCTICA FORMATIVA: INGENIERÍA DE SOFTWARE V} \\
    \fontsize{6.5}{8.5}\selectfont\color{icesiMuted}Período: 202620 $\bullet$ NRC: 11005 $\bullet$ Semana 11 $\bullet$ Sesión 21 $\bullet$ Duración: 2 Horas
  \end{minipage}%
}
\fancyhead[R]{\fontsize{7.5}{9.5}\selectfont\color{icesiMuted}Pág. \thepage\ de \pageref{LastPage}}

% Barra horizontal multicolor
\renewcommand{\headrule}{%
  \vspace{0.08cm}%
  \noindent%
  \textcolor{icesiBlue}{\rule{0.60\textwidth}{1.8pt}}%
  \textcolor{icesiOrange}{\rule{0.18\textwidth}{1.8pt}}%
  \textcolor{icesiGreen}{\rule{0.11\textwidth}{1.8pt}}%
  \textcolor{icesiPurple}{\rule{0.11\textwidth}{1.8pt}}%
}
\renewcommand{\footrulewidth}{0.4pt}
\fancyfoot[L]{\fontsize{7.5}{9.5}\selectfont\color{icesiMuted}Ingeniería de Software V (09776 - TIC) $\bullet$ Taller Formativo: Despliegue Blue-Green}
\fancyfoot[R]{\fontsize{7.5}{9.5}\selectfont\color{icesiMuted}Universidad Icesi $\bullet$ DCSI}

% =========================================================================
% COMANDOS DE SECCIONES CON BADGE ESTILO PROGRAMA ICESI
% =========================================================================
\newcommand{\sectionBadge}[1]{%
  \refstepcounter{section}%
  \vspace{0.16cm}%
  \noindent%
  \colorbox{icesiBlue}{\color{white}\rule[-0.08cm]{0pt}{0.42cm}\hspace{0.16cm}\textbf{\normalsize \thesection}\hspace{0.16cm}}%
  \hspace{0.18cm}%
  {\fontsize{9.5}{12}\selectfont\textbf{\color{icesiDark}#1}}\\[-0.10cm]%
  \noindent{\color{icesiDark}\rule{\textwidth}{0.8pt}}\vspace{0.06cm}%
}

\titleformat{\subsection}
  {\color{icesiBlue}\normalfont\normalsize\bfseries}
  {\thesubsection}{0.4em}{}
\titlespacing*{\subsection}{0pt}{0.4ex plus 0.2ex minus .1ex}{0.2ex plus .1ex}

\newcolumntype{C}[1]{>{\centering\arraybackslash}p{#1}}
\newcolumntype{L}[1]{>{\raggedright\arraybackslash}p{#1}}
\newcolumntype{Y}{>{\centering\arraybackslash}X}

\begin{document}

% =========================================================================
% TÍTULO Y PRESENTACIÓN INICIAL (PÁGINA 1)
% =========================================================================
\noindent
{\fontsize{8.2}{10.5}\selectfont\textbf{\color{icesiBlue}GUÍA PRÁCTICA FORMATIVA: ESTRATEGIA DE DESPLIEGUE BLUE-GREEN (ZERO-DOWNTIME)}}\\[0.04cm]
{\fontsize{13.5}{16}\selectfont\textbf{Ingeniería de Software V --- Sesión 21 (Semana 11)}}\\[0.05cm]
{\fontsize{8}{10}\selectfont\color{icesiMuted}\textbf{Enrutamiento Atómico en Nginx, Pre-Calentamiento, Health Checks y Rollback Instantáneo}}\\[0.02cm]
{\color{icesiRowLine}\hrule height 0.6pt}\vspace{0.04cm}

\noindent
{\color{icesiBlue}\textbf{\rule{2.5pt}{8.0pt}}}~\textbf{\fontsize{8}{10}\selectfont FACULTAD DE INGENIERÍA, DISEÑO Y CIENCIAS APLICADAS} \hfill
{\color{icesiOrange}\textbf{\rule{2.5pt}{8.0pt}}}~\textbf{\fontsize{8}{10}\selectfont DEPARTAMENTO DE COMPUTACIÓN Y SISTEMAS INTELIGENTES}\\[0.02cm]
{\color{icesiRowLine}\hrule height 0.6pt}\vspace{0.04cm}

% =========================================================================
% TABLA DE METADATOS FORMATIVOS
% =========================================================================
\noindent
{\color{icesiBlue}\textbf{\rule{2.5pt}{8.0pt}}}~\textbf{\fontsize{8}{10}\selectfont INFORMACIÓN DE LA ACTIVIDAD} \hfill {\fontsize{7}{8.5}\selectfont\color{icesiMuted}\textbf{IDENTIFICACIÓN}}\\[-0.22cm]
{\color{icesiBorder}\hrule height 0.5pt}\vspace{0.03cm}

\noindent
\fontsize{7.0}{8.8}\selectfont
\begin{tabularx}{\textwidth}{>{\bfseries\color{icesiDark}\fontsize{7.0}{8.8}\selectfont}L{2.5cm}X!{\color{icesiBlue}\vrule width 0.8pt}>{\bfseries\color{icesiDark}\fontsize{7.0}{8.8}\selectfont}L{2.5cm}X}
CÓDIGO: & 09776 - TIC & PERÍODO: & 202620 (Semestre 2026-2) \\ \hline
\rowcolor{icesiLight}
NRC / GRUPO: & 11005 / Grupo 001 & PONDERACIÓN: & \textbf{Formativa} (Sin calificación sumativa) $\bullet$ Duración: 2 Horas \\ \hline
POLÍTICA IAG: & Nivel 3: Colaboración Asistida & MODALIDAD: & Parejas / Laboratorio Práctico Guiado \\ \hline
\rowcolor{icesiLight}
ESTUDIANTE 1: & \rule{4.5cm}{0.4pt} & CÓDIGO: & \rule{2.2cm}{0.4pt} \\ \hline
ESTUDIANTE 2: & \rule{4.5cm}{0.4pt} & CÓDIGO: & \rule{2.2cm}{0.4pt} \\ \hline
\end{tabularx}

\vspace{0.03cm}

% =========================================================================
% SECCIÓN 1: CONTEXTO Y ARQUITECTURA BLUE-GREEN
% =========================================================================
\sectionBadge{Contexto del Problema \& Arquitectura Blue-Green}

\subsection{Descripción del Reto de Despliegue}
Tradicionalmente, actualizar una aplicación implicaba detener el proceso anterior, compilar o desplegar el nuevo binario y reiniciar el servicio. Este enfoque destructivo (\textit{In-Place / Recreate}) ocasiona tiempos de caída inevitables (\textit{Downtime}), desconexión de usuarios y errores HTTP 502 Bad Gateway.

La estrategia \textbf{Blue-Green Deployment} resuelve este problema manteniendo dos entornos de producción arquitectónicamente idénticos:
\begin{itemize}[itemsep=0.2pt,topsep=0.5pt]
  \item \textbf{Entorno Blue (Producción Activa):} Aloja la versión actual (\texttt{v1.0.0}) y recibe el 100\% del tráfico real a través del router de entrada.
  \item \textbf{Entorno Green (Entorno Pasivo / Staging):} Aloja la nueva versión (\texttt{v2.0.0}). Permite realizar pre-calentamiento de cachés, pruebas de humo y validación de endpoints de salud sin alterar a los clientes.
  \item \textbf{Conmutación Atómica de Tráfico:} Una vez validado Green, el enrutador (Nginx) redirige el 100\% de las conexiones a Green en milisegundos mediante una recarga en caliente (\texttt{nginx -s reload}) sin cortar peticiones activas.
  \item \textbf{Rollback Instantáneo:} Si la nueva versión presenta anomalías imprevistas, el router conmutará de vuelta a Blue en menos de 1 segundo.
\end{itemize}

\vspace{0.03cm}
\noindent
\fontsize{7.0}{8.8}\selectfont
\begin{tabularx}{\textwidth}{|L{2.8cm}|X|L{3.5cm}|}
\hline
\rowcolor{icesiLight}
\textbf{Componente} & \textbf{Tecnología \& Contenedor} & \textbf{Puerto / Endpoint} \\ \hline
\textbf{Router / Gateway} & Nginx 1.25 Alpine (Reverse Proxy) & Puerto \texttt{8080} (\texttt{http://localhost:8080}) \\ \hline
\textbf{Archivo Auxiliar} & \texttt{conf.d/upstream.conf} (Dinámico) & Inclusión modular de upstream \\ \hline
\textbf{Servicio Blue} & Microservicio v1.0.0 (\texttt{app-blue}) & Puerto interno \texttt{8081} (\texttt{/api/info}) \\ \hline
\textbf{Servicio Green} & Microservicio v2.0.0 (\texttt{app-green}) & Puerto interno \texttt{8082} (\texttt{/api/info}) \\ \hline
\textbf{Orquestación} & Docker Compose (Red interna \texttt{deploy-net}) & Aislamiento y persistencia \\ \hline
\end{tabularx}

\newpage

% =========================================================================
% PÁGINA 2: PROCEDIMIENTO EXPERIMENTAL PASO A PASO
% =========================================================================
\sectionBadge{Procedimiento Experimental Paso a Paso}

\subsection{Fase 1: Puesta en Marcha del Entorno con Docker Compose (15 min)}
Ubíquese en el directorio \texttt{codigo\_base/} e inspeccione el archivo \texttt{docker-compose.yml}:
\begin{lstlisting}[style=codeStyle]
services:
  nginx-router:
    image: nginx:1.25-alpine
    container_name: nginx-router
    ports:
      - "8080:80"
    volumes:
      - ./nginx/nginx.conf:/etc/nginx/nginx.conf:ro
      - ./nginx/conf.d:/etc/nginx/conf.d
    depends_on:
      - app-blue
      - app-green

  app-blue:
    build: ./app-blue
    container_name: app-blue
    ports:
      - "8081:8080" # Puerto expuesto para verificacion y smoke test

  app-green:
    build: ./app-green
    container_name: app-green
    ports:
      - "8082:8080" # Puerto expuesto para verificacion y smoke test
\end{lstlisting}

Inicie los contenedores y compruebe su estado:
\begin{lstlisting}[style=bashstyle]
docker compose up -d --build
docker compose ps
\end{lstlisting}

\subsection{Fase 2: Arquitectura Modular Nginx y Archivo Auxiliar (15 min)}
\textit{Fundamento de Ingeniería:} Para evitar errores de inodo bloqueado (\textit{Device or resource busy}) al actualizar archivos montados en Docker, Nginx delega la directiva \texttt{upstream} a un archivo auxiliar montado dentro del directorio \texttt{conf.d}:
\begin{lstlisting}[style=codeStyle]
# nginx/nginx.conf: Incluye la definicion del upstream dinamicamente
http {
    include /etc/nginx/conf.d/upstream.conf;
    server {
        listen 80;
        location / {
            proxy_pass http://backend_app;
            proxy_set_header Host $host;
            add_header X-Handled-By $upstream_addr always;
        }
    }
}

# nginx/conf.d/upstream.conf: Archivo auxiliar (inicialmente apunta a Blue)
upstream backend_app {
    server app-blue:8080;
}
\end{lstlisting}

Consulte el router principal a través de \texttt{curl}:
\begin{lstlisting}[style=bashstyle]
curl -i http://localhost:8080/api/info
\end{lstlisting}
Debe observar la cabecera \texttt{X-Handled-By: ...} y la respuesta activa desde el entorno \texttt{BLUE} (versión \texttt{1.0.0}).

\newpage

\subsection{Fase 3: Verificación Previa y Pre-Calentamiento en Green (20 min)}
Antes de transferir a los usuarios, la versión \texttt{v2.0.0} desplegada en Green debe ser sometida a un Smoke Test:
\begin{lstlisting}[style=bashstyle]
# Comprobar salud del microservicio Green antes del swap:
curl -f http://localhost:8082/api/health
curl -s http://localhost:8082/api/info | jq .
\end{lstlisting}
Si el endpoint no responde HTTP 200 con estado \texttt{UP}, la conmutación se cancela inmediatamente, protegiendo a los usuarios.

\subsection{Fase 4: Ejecución del Switch Atómico con Tráfico Sintético (30 min)}
Para evidenciar que la recarga de Nginx ocurre con cero tiempo de inactividad, abra una terminal y ejecute el generador de tráfico en bucle:
\begin{lstlisting}[style=bashstyle]
# Terminal 1: Generador continuo de peticiones
./test_traffic.sh
\end{lstlisting}

En una segunda terminal, ejecute el script de conmutación atómica \texttt{./switch.sh}:
\begin{lstlisting}[style=bashstyle]
#!/usr/bin/env bash
# 1. Comprobacion de salud estricta en Green
echo "[1/3] Verificando salud de entorno Green..."
curl -sf http://localhost:8082/api/health || { echo "[ERROR] Green no saludable"; exit 1; }

# 2. Modificacion atomica del archivo auxiliar (sin bloquear inodos)
echo "[2/3] Actualizando archivo auxiliar de upstream hacia app-green..."
sed -i 's/app-blue:8080/app-green:8080/' ./nginx/conf.d/upstream.conf

# 3. Recarga atomica en caliente sin cortar conexiones existentes
echo "[3/3] Recargando Nginx en caliente (Zero Downtime)..."
docker exec nginx-router nginx -s reload
echo "[OK] Conmutacion completada. Produccion ahora corre sobre GREEN (v2.0.0)"
\end{lstlisting}
Observe la Terminal 1: las peticiones cambian de \texttt{BLUE} a \texttt{GREEN} instantáneamente sin que ninguna arroje error HTTP 502 ni tiempo de espera excesivo.

\subsection{Fase 5: Simulación de Incidente y Rollback Instantáneo (20 min)}
Simule un fallo en caliente ejecutando el script de reversión:
\begin{lstlisting}[style=bashstyle]
./rollback.sh
\end{lstlisting}
El script revierte el archivo auxiliar \texttt{upstream.conf} hacia \texttt{app-blue:8080} y recarga Nginx en menos de 1 segundo.

\newpage

% =========================================================================
% PÁGINA 4: CUESTIONARIO Y CHECKLIST FORMATIVO
% =========================================================================
\sectionBadge{Cuestionario de Reflexión Conceptual \& Checklist Formativo}

\subsection{Preguntas de Reflexión Técnica}
\begin{enumerate}[leftmargin=*,itemsep=2.5pt]
  \item \textbf{Costo de Infraestructura:} ¿Qué implicaciones presupuestales tiene la duplicación al 100\% de recursos que exige Blue-Green en clústeres de cientos de microservicios? ¿Cómo mitiga la nube elástica este costo?
  \item \textbf{Persistencia de Sesiones:} Si los usuarios almacenan carritos de compra o sesiones en la memoria RAM del contenedor Blue, ¿qué ocurriría al conmutar a Green? ¿Por qué la arquitectura desacoplada Twelve-Factor (Factor VI) exige sesiones en servicios externos (ej. Redis)?
  \item \textbf{Patrón Expand \& Contract en BD:} Si la versión Green requiere una columna nueva en la base de datos relacional compartida, ¿por qué no se debe renombrar ni eliminar una columna previa mientras Blue continúe activo?
  \item \textbf{Arquitectura de Recarga e Inodos en Docker:} ¿Por qué ejecutar \texttt{sed -i} sobre un archivo montado puntualmente (\textit{single file bind mount}) genera el error \textit{Device or resource busy}? ¿Cómo resuelve este problema el patrón arquitectónico de archivo auxiliar incluido (\texttt{upstream.conf}) dentro de un directorio montado (\texttt{conf.d})?
\end{enumerate}

\vspace{0.05cm}
\subsection{Checklist de Logro Formativo (Autoevaluación)}
\noindent
\small
\begin{tabularx}{\textwidth}{|L{4.2cm}|X|C{2.5cm}|}
\hline
\rowcolor{icesiLight}
\textbf{Objetivo Formativo} & \textbf{Evidencia Observable} & \textbf{Estado} \\ \hline
\textbf{1. Orquestación Multi-Entorno} & Contenedores Nginx, Blue (v1) y Green (v2) corriendo en red aislada. & [ ] Logrado \\ \hline
\textbf{2. Archivo Auxiliar Upstream} & Uso de \texttt{conf.d/upstream.conf} modular sin errores de inodo bloqueado. & [ ] Logrado \\ \hline
\textbf{3. Smoke Test en Staging} & Validación de endpoint \texttt{/api/health} en Green antes del switch. & [ ] Logrado \\ \hline
\textbf{4. Zero Downtime Swap} & Conmutación con \texttt{nginx -s reload} sin pérdida de paquetes HTTP. & [ ] Logrado \\ \hline
\textbf{5. Rollback Instantáneo} & Retorno a Blue en $<1$ segundo mediante reversión de archivo auxiliar. & [ ] Logrado \\ \hline
\end{tabularx}

\end{document}
